%% ---------------------------------------------------------------------
\section{Por qué retroalimentar: sensitividad y error}
\label{sec:er-sensitividad}
%% ---------------------------------------------------------------------

Analizar un sistema de control es evaluar tres cosas: la estabilidad (sin ella las otras dos no tienen sentido), la respuesta transitoria y el \textbf{error en estado estacionario}. Las dos primeras ocuparon los capítulos~\ref{cap:po}, \ref{cap:so} y~\ref{cap:es}; este capítulo trata la tercera.

En malla abierta nada corrige una perturbación ni un cambio en los parámetros del proceso. La retroalimentación sí: reduce la sensibilidad a variaciones de la planta, mejora el rechazo de perturbaciones, reduce el error en estado estacionario y facilita ajustar el transitorio. A cambio, hace llegar al lazo el ruido del sensor. La \cref{fig:er-rdn} muestra las tres señales exógenas: la referencia $R$, la perturbación $D$ y el ruido de medición $N$. Definimos el \clave{error de seguimiento} como
\begin{equation}
  E(s)=R(s)-Y(s),
  \label{eq:er-error}
\end{equation}
lo que nos interesa del desempeño. Con $\Hs=1$, por superposición,
\begin{equation}
  E=\underbrace{\frac{1}{1+\Gc\Gp}}_{\Sens(s)}R-\frac{\Gp}{1+\Gc\Gp}D+\underbrace{\frac{\Gc\Gp}{1+\Gc\Gp}}_{\Tlc(s)}N
  =\Sens R-\Sens\Gp D+\Tlc N .
  \label{eq:er-sensitividad}
\end{equation}

\begin{figure}[ht]
\centering
\includegraphics[width=0.8\linewidth]{cap06/lazo_rdn}
\caption{Lazo con referencia $R$, perturbación $D$ a la entrada del proceso y ruido de medición $N$.}
\label{fig:er-rdn}
\end{figure}

\begin{definicion}[label=def:er-sensitividad]{Sensitividad y sensitividad complementaria}
Con retroalimentación unitaria, la \textbf{sensitividad} es $\Sens(s)=\dfrac{1}{1+\Gc\Gp}$ y la \textbf{sensitividad complementaria} es $\Tlc(s)=\dfrac{\Gc\Gp}{1+\Gc\Gp}$, que es la propia función de lazo cerrado $Y/R$. Para todo $s$,
\[
  \Sens(s)+\Tlc(s)=1 .
\]
\end{definicion}

Todos los términos de \eqref{eq:er-sensitividad} tienen el mismo denominador: la ecuación característica $1+\Gc\Gp=0$ es la misma para cualquier entrada. Y la identidad $\Sens+\Tlc=1$ dice que no se pueden hacer pequeñas las dos a la misma frecuencia. El compromiso se resuelve separando frecuencias: $\Gc\Gp$ grande donde viven las perturbaciones y las referencias (bajas frecuencias), para que $\Sens$ sea pequeña; $\Gc\Gp$ pequeña donde vive el ruido (altas frecuencias), para que $\Tlc$ lo sea.

\begin{ejemplo}[label=ej:er-sens]{Dos ganancias}
Con $\Gp=1/((s+1)(s+10))$ y $\Gc=K$, compare $K=10$ y $K=200$ ante una perturbación $d=0.5\sin(0.2t)$ y un ruido $n=0.01\sin(50t)$, con referencia constante.
\solucion La amplitud de cada aporte es la amplitud de la señal por la magnitud de su función a esa frecuencia. Con $K=200$: perturbación $0.5\,|\Sens\Gp|(0.2)=0.0024$ y ruido $0.01\,|\Tlc|(50)=0.00085$. Con $K=10$: $0.025$ y $0.00004$. La ganancia alta rechaza diez veces mejor la perturbación y deja pasar veinte veces más ruido (\cref{fig:er-sens}); aquí gana, porque el ruido es pequeño. Si su amplitud pasara de $0.28$ (sumando los aportes como cota), convendría $K=10$.
\end{ejemplo}

\begin{figure}[ht]
\centering
\includegraphics[width=0.85\linewidth]{cap06/f_sensibilidad}
\caption{Ejemplo~\ref{ej:er-sens}: $|\Sens|$ y $|\Tlc|$ para $K=10$ y $K=200$. Donde una es pequeña, la otra vale casi $1$.}
\label{fig:er-sens}
\end{figure}

\codigorepo{cap06_error/C6_Ej6_Sensibilidad.m}{aportes}{ejemplo~\ref{ej:er-sens}}

La sensitividad y la robustez se estudian a fondo en cursos posteriores. Aquí nos quedamos con su consecuencia más directa: el error que queda cuando todo se asienta. En lo que sigue (salvo en la \cref{sec:er-perturbacion}) el lazo no tiene ni perturbación ni ruido.
